
\chapter{CÁC NGHIÊN CỨU LIÊN QUAN VÀ PHÂN TÍCH KHOẢNG TRỐNG}
\label{chap:related_works}

Sự hội tụ giữa Robot học, Thực tế ảo (VR), và Bản sao số (Digital Twin) đã thúc đẩy hàng loạt nghiên cứu trong thập kỷ qua. Chương này phân tích sâu các tài liệu nghiên cứu hiện có theo bốn nhóm chủ đề: (1) Giao diện điều khiển, (2) Tích hợp Bản sao số, (3) Tối ưu Middleware và quản lý độ trễ, (4) Haptic Feedback. Mỗi nghiên cứu được đánh giá về điểm mạnh và hạn chế cụ thể để làm rõ vị trí và hướng tiếp cận của luận văn.

\section{Giao diện điều khiển robot từ xa}
\label{sec:related_interface}

Giao diện điều khiển (Control Interface) là yếu tố quyết định trải nghiệm của người vận hành và hiệu quả của hệ thống teleoperation.

Chen et al. (2023) \cite{chen2023comparing} thực hiện so sánh trực tiếp ba phương thức điều khiển trên cánh tay JAKA Minicobo: Giao diện 2D (GUI) của nhà sản xuất, Cử chỉ tay (Hand Gestures), và Bộ điều khiển VR (VR Controllers) với kính Meta Quest 2. Thử nghiệm với 16 người dùng trên nhiệm vụ cụ thể (điều khiển cánh tay câu cá) cho thấy VR Controllers đạt thời gian hoàn thành ngắn hơn đáng kể so với GUI 2D, và ổn định hơn so với Cử chỉ tay do giảm được nhiễu bàn tay tự nhiên. \textit{Điểm mạnh}: nghiên cứu có dữ liệu thực nghiệm với nhóm người dùng đa dạng, phương pháp đo lường khách quan. \textit{Hạn chế}: sử dụng Meta Quest 2 (chip thế hệ cũ, hiệu năng thấp hơn Quest 3) và ROS1, nhiệm vụ thử nghiệm đơn giản và không đại diện cho môi trường công nghiệp thực.

Hetrick et al. (2020) \cite{hetrick2020comparing} nghiên cứu sự khác biệt giữa hai phương thức điều khiển trên robot Baxter hai tay với kính HTC Vive: điều khiển theo quỹ đạo (trajectory-based) và điều khiển theo vị trí (goal-based). Kết quả cho thấy không có sự khác biệt lớn về hiệu suất, nhưng người có kinh nghiệm thiên về phương thức goal-based do tốc độ nhanh hơn. \textit{Điểm mạnh}: phân tích được ảnh hưởng của kinh nghiệm người dùng lên cognitive load và workload. \textit{Hạn chế}: robot Baxter hai tay không phổ biến trong bối cảnh cánh tay công nghiệp đơn, HTC Vive là thiết bị đã lỗi thời.

Su et al. (2022) \cite{su2022mixed} đề xuất hệ thống Mixed Reality kết hợp imitation-based mapping, ánh xạ cử chỉ tay người điều khiển sang chuyển động cánh tay robot một cách tự nhiên. \textit{Điểm mạnh}: tính trực quan cao khi ánh xạ trực tiếp từ cử chỉ người. \textit{Hạn chế}: độ trễ cao do xử lý MR và phụ thuộc vào PC đồ họa mạnh, không khả thi cho triển khai standalone.

Zhang et al. (2023) \cite{zhang2023research} giới thiệu hệ thống VR 5 chiều (5D) cho teleoperation với tích hợp đa modal bao gồm cả phản hồi lực. \textit{Điểm mạnh}: cung cấp trải nghiệm toàn diện với nhiều kênh phản hồi. \textit{Hạn chế}: phần cứng đặc thù, không mở nguồn (not open-source), khó tái tạo trong môi trường thực nghiệm độc lập.

\textbf{Nhận xét:} Kết quả tổng hợp từ các nghiên cứu trên cho thấy VR Controllers mang lại kết quả tốt nhất về tốc độ và ổn định so với cử chỉ tay và GUI 2D. Đây là lý do luận văn lựa chọn Meta Quest 3 với VR Controllers làm giao diện điều khiển chính, thay vì hand tracking.

\section{Tích hợp Bản sao số và nhận thức tình huống}
\label{sec:related_digital_twin}

Bản sao số (Digital Twin) là thành phần cốt lõi để tăng nhận thức tình huống (situational awareness) cho người vận hành trong môi trường VR.

Naceri et al. (2021) \cite{naceri2021vicarios} đề xuất \textit{Vicarios Interface} trên kính HTC Vive Pro, cho phép người điều khiển ``dịch chuyển'' (teleport) trong không gian ảo để thay đổi góc nhìn tự do và render môi trường bằng Point Cloud từ nhiều camera. \textit{Điểm mạnh}: giải quyết bài toán điểm mù (occlusion) hiệu quả, người dùng đánh giá cao tính trực quan của góc nhìn 3D tự do. \textit{Hạn chế}: sử dụng ROS1 với ROSBridge gây bottleneck băng thông khi truyền dữ liệu Point Cloud khối lượng lớn, đòi hỏi PC đồ họa mạnh để render real-time.

Huang et al. (2020) \cite{multi-view} xây dựng nền tảng AR đa camera với ZED stereo camera (toàn cục) và RealSense (cục bộ gắn tay gắp), sử dụng Point Cloud Depth Projection để tạo cảm giác ``nhìn xuyên qua'' khi tay gắp bị che khuất. \textit{Điểm mạnh}: Picture-in-Picture kết hợp Depth Projection cải thiện nhận thức tình huống trong môi trường có vật cản. \textit{Hạn chế}: thiết kế Picture-in-Picture khi không hợp lý làm tăng cognitive load thay vì giảm, như kết quả nghiên cứu cũng chỉ ra.

Torrejón et al. (2026) \cite{torrejon2026teleoperation} xây dựng hệ thống teleoperation tay đôi với Point Cloud mật độ cao ($>$1 triệu điểm) được render bằng Unreal Engine 5.6, sử dụng kính Meta Quest 3. \textit{Điểm mạnh}: tái tạo môi trường vật lý với mức độ chân thực cao, cung cấp nhận thức tình huống phong phú. \textit{Hạn chế}: yêu cầu PC có GPU NVIDIA RTX 3070 trở lên, độ trễ hệ thống đạt 190ms (sát ngưỡng thoải mái cho VR), không thể chạy standalone.

Kawshan và Peng (2026) \cite{kawshan2026digital} xây dựng Digital Twin trên Unity kết nối với JetCobot 7-DoF qua ROS1. Nghiên cứu cung cấp benchmark chi tiết về độ trễ end-to-end đạt 144ms và xác nhận cơ chế Master-Slave của ROS1 là bottleneck. \textit{Điểm mạnh}: đánh giá độ trễ đầy đủ, cung cấp số liệu tham chiếu quan trọng. \textit{Hạn chế}: bottleneck ROS1 đã được xác nhận nhưng chưa được giải quyết trong phạm vi nghiên cứu này.

Singh et al. (2025) \cite{singh2025comparative} so sánh Digital Twin trong Unity và Gazebo về simulation fidelity. \textit{Điểm mạnh}: phân tích học thuật có hệ thống về độ trung thực mô phỏng. \textit{Hạn chế}: chỉ dừng ở mô phỏng, không triển khai vòng điều khiển kín (closed-loop) với phần cứng thực tế.

Ali et al. (2025) \cite{ali2025digital} sử dụng ROS2 và Digital Twin để huấn luyện Reinforcement Learning cho robot sản xuất. \textit{Điểm mạnh}: tiên phong áp dụng ROS2 DDS cho Digital Twin, cho thấy ưu thế so với ROS1. \textit{Hạn chế}: hoạt động ở chế độ offline (huấn luyện AI), không phục vụ người vận hành tương tác VR trực tiếp.

\textbf{Nhận xét:} Luận văn này sử dụng cách tiếp cận Digital Twin đơn giản và thực tế hơn: mô hình URDF/CAD được load sẵn trong Meta Quest 3, đồng bộ trạng thái qua Joint States và camera stream qua WebSocket. Cách tiếp cận này không cung cấp độ phong phú môi trường như Point Cloud (so với Naceri, Torrejón), nhưng cho phép chạy standalone trên Quest 3 đạt 66,1 FPS với 60\% tải GPU mà không cần PC đồ họa --- ưu thế trực tiếp về tính khả thi và chi phí triển khai.

\section{Tối ưu hóa Middleware và Quản lý Độ trễ}
\label{sec:related_middleware}

Việc lựa chọn nền tảng giao tiếp quyết định tính thời gian thực của toàn hệ thống.

Kawshan và Peng (2026) \cite{kawshan2026digital} đã định lượng được bottleneck của ROS1: tổng độ trễ end-to-end đạt 144ms trong hệ thống Unity--ROS1. Cơ chế Master-Slave TCP/IP của ROS1 bộc lộ điểm nghẽn nặng khi phải truyền đồng thời kinematics tần số cao và dữ liệu hình ảnh.

Luu et al. (2025) \cite{luu2025enhancing} phân tích hệ thống teleoperation trong Industrial IoT, nhấn mạnh ảnh hưởng của hạ tầng mạng lên độ trễ điều khiển. \textit{Điểm mạnh}: xem xét vấn đề từ góc độ IoT infrastructure toàn diện. \textit{Hạn chế}: không sử dụng hardware interface chuẩn hóa như ros2-control.

Ali et al. (2025) \cite{ali2025digital} áp dụng ROS2 DDS cho Digital Twin, đặt tiền đề kỹ thuật cho việc thay thế ROS1. Macenski et al. (2022) \cite{doi:10.1126/scirobotics.abm6074} trình bày kiến trúc và thiết kế của ROS2, khẳng định DDS là cơ sở cho giao tiếp phân tán, QoS thời gian thực, và tách biệt luồng dữ liệu.

Li et al. (2025) \cite{reality-fusion} đề xuất giải pháp truyền phát môi trường dưới dạng 3D Gaussian Splats (3DGS) cho robot di động, cải thiện chất lượng hình ảnh. \textit{Điểm mạnh}: hình ảnh môi trường chân thực với băng thông tối ưu hơn Point Cloud. \textit{Hạn chế}: 3DGS được thiết kế cho robot di động, không phù hợp với cánh tay công nghiệp cố định; đòi hỏi bước tái tạo 3DGS từ camera trước khi triển khai.

\textbf{Nhận xét:} Luận văn lựa chọn ROS2 (DDS) làm nền tảng để khắc phục trực tiếp bottleneck của ROS1. Nhờ DDS, luồng camera stream (băng thông cao) và luồng Joint States (tần số cao) được xử lý độc lập qua các topic riêng biệt với QoS phù hợp, đảm bảo tính thời gian thực cho điều khiển mà không bị ảnh hưởng bởi tải hình ảnh.

\section{Haptic Feedback và Phản hồi cảm giác}
\label{sec:related_haptic}

Phản hồi xúc giác (Haptic Feedback) là hướng nghiên cứu bổ sung quan trọng trong teleoperation, giúp người vận hành cảm nhận được lực tiếp xúc từ xa.

Xu et al. (2025) \cite{xu2025immersive} phát triển hệ thống VR bimanual (hai tay) tích hợp với găng tay Dexmo để truyền phản hồi lực từ cảm biến lực trên robot. \textit{Điểm mạnh}: phản hồi xúc giác thực sự giúp tăng độ chính xác trong tác vụ lắp ráp tinh vi, người dùng cảm nhận được khi robot chạm vật thể. \textit{Hạn chế}: thiết bị Dexmo có chi phí cao, cần cảm biến lực/moment (F/T sensor) gắn trực tiếp trên robot, thêm độ trễ xử lý tín hiệu haptic.

Papakonstantinou et al. (2025) \cite{papakonstantinou2025effect} phân tích định lượng tác động của xúc giác lên hiệu suất điều khiển cánh tay đơn. \textit{Điểm mạnh}: số liệu định lượng rõ ràng về mức độ cải thiện khi có haptic. \textit{Hạn chế}: đòi hỏi F/T sensor gắn trực tiếp vào robot, làm tăng độ phức tạp tích hợp.

MIT ARCLab (2025) \cite{arclab2025beavr} giới thiệu BEAVR: hệ thống teleoperation VR bimanual mã nguồn mở, hỗ trợ đa nền tảng. \textit{Điểm mạnh}: open-source, khả năng tái sử dụng cao. \textit{Hạn chế}: được thiết kế cho robot nghiên cứu (Franka Emika) vốn có API real-time sẵn có, không áp dụng trực tiếp cho robot công nghiệp đời cũ với bộ điều khiển hạn chế.

\textbf{Nhận xét:} Cánh tay Denso VS-6577/RC5 không được trang bị F/T sensor tinh vi, và chi phí tích hợp haptic là rào cản thực tế trong bối cảnh của luận văn. Thay vào đó, luận văn tập trung tối ưu hóa phản hồi thị giác (visual feedback) thông qua Digital Twin thời gian thực và camera stream --- dùng thị giác bù đắp xúc giác, phù hợp với điều kiện phần cứng và kinh tế của đề tài.

\section{Bảng so sánh các nghiên cứu liên quan}
\label{sec:related_table}

Bảng \ref{tab:literature_review} tổng hợp và đối chiếu các nghiên cứu tiêu biểu với hệ thống được đề xuất trong luận văn.

\begin{table}[H]
\centering
\caption{So sánh các nghiên cứu liên quan về VR Teleoperation với luận văn hiện tại}
\label{tab:literature_review}
\resizebox{\textwidth}{!}{%
\begin{tabular}{|p{3.2cm}|p{2.5cm}|p{2.5cm}|p{3cm}|p{3.8cm}|p{3.8cm}|}
\hline
\textbf{Tác giả / Năm} & \textbf{Nền tảng Robot} & \textbf{Middleware \& Giao tiếp} & \textbf{Thiết bị VR/AR} & \textbf{Điểm mạnh} & \textbf{Hạn chế / Khác biệt với Luận văn} \\ \hline

Chen et al. (2023) \cite{chen2023comparing} & JAKA Minicobo & ROS1, UDP & Meta Quest 2 & Dữ liệu thực nghiệm 16 người dùng, chứng minh VR Controller hiệu quả hơn GUI 2D. & Dùng ROS1 và thiết bị VR cũ; nhiệm vụ thử nghiệm đơn giản. Ours: ROS2 + Quest 3. \\ \hline

Naceri et al. (2021) \cite{naceri2021vicarios} & UR5, SoftHand & ROS1, ROSBridge, Unreal & HTC Vive Pro & Teleporting góc nhìn tự do; render Point Cloud cải thiện nhận thức tình huống. & ROS1/ROSBridge gây bottleneck băng thông; cần PC đồ họa mạnh. Ours: ROS2 DDS, standalone Quest 3. \\ \hline

Torrejón et al. (2026) \cite{torrejon2026teleoperation} & P3bot (tay đôi) & RoboComp, Unreal 5.6 & Meta Quest 3 & Point Cloud $>$1M điểm, môi trường 3D chân thực. & Cần GPU RTX 3070+, độ trễ 190ms, không standalone. Ours: URDF lightweight, 66,1 FPS standalone. \\ \hline

Xu et al. (2025) \cite{xu2025immersive} & Robot tay đôi, Dexmo & Unity, custom protocol & HTC Vive & Haptic Feedback thực sự tăng độ chính xác lắp ráp tinh vi. & Cần F/T sensor và găng tay đắt tiền. Ours: tối ưu visual feedback thay cho haptic. \\ \hline

Kawshan \& Peng (2026) \cite{kawshan2026digital} & JetCobot (7-DoF) & ROS1, ROS-TCP, Unity & Desktop/VR & Benchmark độ trễ chi tiết (144ms end-to-end). & Bottleneck ROS1 được xác nhận nhưng chưa giải quyết. Ours: chuyển hoàn toàn sang ROS2 DDS. \\ \hline

Luu et al. (2025) \cite{luu2025enhancing} & Robot công nghiệp & ROS, IoT network & -- & Phân tích ảnh hưởng hạ tầng mạng IoT lên độ trễ. & Không dùng hardware interface chuẩn hóa. Ours: ros2-control + MoveIt2. \\ \hline

\textbf{Luận văn (Ours)} & \textbf{Denso VS-6577 \& RC5} & \textbf{ROS2 (DDS), ROS-Sharp, UART} & \textbf{Meta Quest 3, Unity 6} & \textbf{Full-stack từ VR đến bộ điều khiển công nghiệp đời cũ; standalone 66,1 FPS; kiến trúc ros2-control chuẩn hóa.} & \textbf{Chưa tích hợp haptic; Trend Learning là giải pháp tạm thời cho vấn đề đồng bộ RC5.} \\ \hline

\end{tabular}%
}
\end{table}

\section{Tổng kết: Khoảng trống nghiên cứu và định hướng của luận văn}
\label{sec:related_gap}

Tổng hợp từ phân tích các nghiên cứu trên, có thể xác định ba khoảng trống nghiên cứu (research gap) mà luận văn hướng đến:

\begin{enumerate}

    \item \textbf{Phần lớn hệ thống teleoperation vẫn dùng ROS1 với bottleneck đã được xác nhận.} Kawshan \& Peng (2026) đã định lượng bottleneck này là 144ms, nhưng không giải quyết. Các công trình dùng ROS2 như Ali et al. (2025) chủ yếu phục vụ huấn luyện AI offline, chưa có công trình nào kết hợp ros2-control với giao diện VR tương tác trực tiếp cho robot công nghiệp. Luận văn lấp đầy khoảng trống này bằng kiến trúc ROS2 (DDS) phân tầng chuẩn hóa.

    \item \textbf{Digital Twin phụ thuộc Point Cloud nặng và PC đồ họa chuyên dụng.} Naceri (2021) và Torrejón (2026) cung cấp nhận thức tình huống phong phú nhưng không khả thi cho triển khai standalone. Luận văn chọn cách tiếp cận URDF lightweight + WebSocket, đạt được tính khả thi cao với 66,1 FPS trên Quest 3 mà không cần PC đồ họa đi kèm.

    \item \textbf{Chưa có công trình nào thực hiện tích hợp full-stack từ VR hiện đại đến bộ điều khiển công nghiệp đời cũ không hỗ trợ real-time API.} Các robot trong nghiên cứu học thuật (UR5, Franka Emika, JAKA) đều có API điều khiển linh hoạt. Denso VS-6577 với RC5 chỉ hỗ trợ giao tiếp ngoại vi qua UART 115200 baud --- một giới hạn phần cứng thực tế mà luận văn giải quyết bằng giao thức UART tùy chỉnh, Ring Buffer, và Trend Learning. Cần thẳng thắn nhìn nhận rằng Trend Learning là giải pháp tạm thời giúp cải thiện đáng kể tính mượt mà của quỹ đạo, nhưng vấn đề đồng bộ giữa tần số MoveIt Servo và tốc độ đáp ứng RC5 vẫn là hướng nghiên cứu mở cần tiếp tục.

\end{enumerate}
